\documentclass[aspectratio=169,11pt]{beamer}

\usepackage[T1]{fontenc}
\usepackage[utf8]{inputenc}
\usepackage[italian]{babel}
\usepackage{lmodern}
\usepackage{booktabs}
\usepackage{tikz}
\usetikzlibrary{positioning,arrows.meta,calc}

\usetheme{Madrid}
\setbeamertemplate{navigation symbols}{}

% footline a una sola casella, colore uniforme
% autori a sinistra, pagina a destra
\setbeamertemplate{footline}{%
  \leavevmode%
  \hbox{%
  \begin{beamercolorbox}[wd=\paperwidth,ht=2.25ex,dp=1ex,leftskip=2ex,rightskip=2ex]{author in head/foot}%
    \usebeamerfont{author in head/foot}\insertshortauthor%
    \hfill%
    \usebeamerfont{page number in head/foot}\usebeamertemplate{page number in head/foot}%
  \end{beamercolorbox}}%
  \vskip0pt%
}

\definecolor{granata}{HTML}{8A2836}
\definecolor{acqua}{HTML}{10585D}
\definecolor{grigio}{HTML}{5C636C}
\usecolortheme[named=granata]{structure}

\tikzset{
  stage/.style={draw=granata!55, fill=granata!7, rounded corners=2pt,
                inner sep=3pt, minimum height=7.5mm, text width=17mm,
                align=center, font=\scriptsize\bfseries},
  nodo/.style={draw=acqua!60, fill=acqua!7, rounded corners=2pt,
               inner sep=3pt, minimum height=7mm, text width=20mm,
               align=center, font=\scriptsize\bfseries},
  term/.style={draw=grigio!60, fill=black!4, rounded corners=2pt,
               inner sep=3pt, minimum height=7mm, text width=16mm,
               align=center, font=\scriptsize},
  sub/.style={font=\tiny, text=grigio, align=center, text width=19mm},
  nota/.style={font=\tiny\itshape, text=acqua, align=center},
  fl/.style={-{Stealth[length=4pt]}, draw=grigio, semithick},
  loop/.style={-{Stealth[length=4pt]}, draw=acqua, semithick, dashed},
  scrivi/.style={-{Stealth[length=4pt]}, draw=granata, semithick, dotted},
}

\title[Sapienza Agents]{Sapienza Agents}
\subtitle{Un sistema multi-agente per domande in linguaggio naturale}
\author[Alessandro Chitarrini, Matteo Crugliano, Davide Gaglione]{%
  \texorpdfstring{%
  \begin{tabular}{@{}ccc@{}}
    Alessandro Chitarrini & Matteo Crugliano & Davide Gaglione \\[0.15em]
    {\tiny\texttt{chitarrini.2129549@studenti.uniroma1.it}} &
    {\tiny\texttt{crugliano.2127105@studenti.uniroma1.it}} &
    {\tiny\texttt{gaglione.2150498@studenti.uniroma1.it}}
  \end{tabular}%
  }{Alessandro Chitarrini, Matteo Crugliano, Davide Gaglione}%
}
\institute[]{Laboratorio di Ingegneria Informatica \\ Sapienza Università di Roma}
\date{2 settembre 2026}

\begin{document}

% ---------------------------------------------------------------- titolo
\begin{frame}[plain]
  \titlepage
\end{frame}

% ---------------------------------------------------------------- orchestratore
\begin{frame}{Orchestratore}
\centering

\begin{tikzpicture}

\node[term] (dom) at (0,0) {Domanda};
\node[nodo] (sup) at (3.0,0) {Supervisor};
\node[nodo] (kg) at (6.3,1.15) {Agente kg};
\node[nodo] (pl) at (6.3,0) {Agente planner};
\node[nodo] (ma) at (6.3,-1.15) {Agente multiapi};
\node[nodo] (syn) at (9.5,0) {Synthesizer};
\node[term] (out) at (12.1,0) {Risposta};

\draw[fl] (dom) -- (sup);
\draw[fl] (sup) -- (kg);
\draw[fl] (sup) -- (pl);
\draw[fl] (sup) -- (ma);
\draw[fl] (kg) -- (syn);
\draw[fl] (pl) -- (syn);
\draw[fl] (ma) -- (syn);
\draw[fl] (syn) -- (out);

\end{tikzpicture}

\end{frame}

% ---------------------------------------------------------------- agente kg
\begin{frame}{Agente knowledge graph}
\centering
\begin{tikzpicture}[node distance=3.2mm]

\node[stage] (cache) {Cache\\semantica};
\node[stage, right=of cache] (link) {Entity\\linking};
\node[stage, right=of link] (schema) {Schema\\retrieval};
\node[stage, right=of schema] (trad) {Traduzione\\text2query};
\node[stage, right=of trad] (exec) {Esecuzione};
\node[stage, right=of exec] (ground) {Grounding};

\foreach \a/\b in {cache/link, link/schema, schema/trad, trad/exec, exec/ground}
  \draw[fl] (\a) -- (\b);

\draw[loop] ([xshift=4mm]exec.north) .. controls +(0,17mm) and +(0,17mm) .. ([xshift=-4mm]trad.north)
  node[midway, above=0.6mm, font=\tiny, text=acqua, fill=white, inner sep=1pt]
  {zero righe $\rightarrow$ rigenerazione ReAct};

\draw[loop] ([xshift=-2mm]exec.north) .. controls +(0,8mm) and +(0,8mm) .. ([xshift=2mm]trad.north)
  node[midway, above=0.6mm, font=\tiny, text=acqua, fill=white, inner sep=1pt]
  {errore $\rightarrow$ riscrivi la query};

\coordinate (dx) at ([xshift=6mm]ground.east);
\coordinate (sx) at ([xshift=-6mm]cache.west);
\coordinate (giu) at ([yshift=-6mm]cache.south);
\draw[scrivi, rounded corners=3pt]
  (ground.east) -- (dx) -- (dx |- giu) -- (sx |- giu) -- (sx) -- (cache.west);
\node[font=\tiny, text=granata, below=0.4mm] at ($(sx |- giu)!0.5!(dx |- giu)$)
  {scrittura in cache, solo se ci sono risultati};

\end{tikzpicture}

\vspace{1.2em}

\begin{columns}[t,onlytextwidth]
\begin{column}{0.44\textwidth}
\scriptsize
\vspace{-5.1em}
\begin{itemize}\setlength\itemsep{1pt}
  \item tre grafi (Wikidata, DBpedia, Neo4j) supportati
  \item zero-shot: nessun fine-tuning
  \item QALD-10: macro-F1 \textbf{0.377} contro \textbf{0.284} senza grafo
  \item lavora in inglese per costruzione
\end{itemize}
\end{column}
\begin{column}{0.52\textwidth}
\centering\scriptsize\renewcommand{\arraystretch}{0.95}
\begin{tabular}{@{}ll@{}}
\toprule
\textbf{modello} & \textbf{dove} \\
\midrule
all-MiniLM-L6-v2 & cache semantica \\
GLiNER & estrazione delle menzioni \\
bge-small-en-v1.5 & schema retrieval, candidati \\
qwen2.5:7b-instruct & filtro e disambiguazione entità \\
qwen2.5-coder:7b & generazione e correzione query \\
\bottomrule
\end{tabular}
\end{column}
\end{columns}
\end{frame}

% ----------------------------------------------------------------  multiapi
\begin{frame}{Agente multiapi}
\centering
todo
\end{frame}

% ---------------------------------------------------------------- planner
\begin{frame}{Agente planner}
\centering
todo
\end{frame}

\end{document}
